\documentclass[10pt,a4paper]{amsart}
\usepackage[margin=19mm]{geometry}
\usepackage{amsmath,amssymb,mathtools}
\usepackage[scr=rsfs,cal=euler]{mathalfa}
\usepackage{xfrac}
\usepackage[unicode,hidelinks,bookmarks=false]{hyperref}
\newcommand{\ie}{i.e.\ }
\newcommand{\cf}{cf.\ }
\newcommand{\resp}{respectively\ }
\newcommand{\Subsection}{\subsection}
\providecommand{\nobreakdash}{\nobreak\hbox{-}}
\setlength{\emergencystretch}{2em}
\multlinegap0pt
\usepackage{float}
\newtheorem{lemma}{Lemma}
\newtheorem{proposition}{Proposition}
\title{Twisted Frobenius-Schur Indicators and Character Degree Sums in Dihedral Groups}
\author{Venkata Subbaiah Yerrapati, Rahul Dixit, Ajay Kumar Shukla}
\date{Source: arXiv:2605.22127v1 (CC BY 4.0)}
\begin{document}
\maketitle

\noindent\emph{Source and licence.} Twisted Frobenius-Schur Indicators and Character Degree Sums in Dihedral Groups. Version arXiv:2605.22127v1 (CC BY 4.0), distributed by arXiv under the Creative Commons Attribution 4.0 International licence, http://creativecommons.org/licenses/by/4.0/, as recorded in the arXiv licence metadata for this exact version and shown on the abstract page https://arxiv.org/abs/2605.22127v1. Full licence notice: \texttt{licenses/arxiv-2605.22127-CC-BY-4.0.txt}.\par\smallskip

\noindent\textit{Editorial scope.} Counted unit: the article's lemma identity\_greater\_other (source0.tex lines 249--251). The article's complete original proof of it is delivered with it (source0.tex lines 252--268, one \texttt{proof} environment of 17 lines). No mathematics is written, repaired or re-derived: the statement and the proof are the article's own lines, in the article's own notation, and no retained line is substituted or re-flowed.  Retained uncounted as the source-local context the proof needs: lines 128--130; lines 167--169; lines 188--197.  Nothing was omitted from any retained span, so the package carries no omission record.  No retained line contains a citation, so the wrapper carries no bibliography and no bibliography entry.  No retained line contains a figure environment, an image-inclusion command or drawing code, so no image is delivered and the package declares no asset dependency.  Editorial formatting only, in addition: an AMS \texttt{amsart} preamble that restates the article's own declarations and macros used by the retained lines, namely the environments \texttt{lemma}, \texttt{proposition} on separate counters, together with the standard packages those lines need.  The retained environments print in this document as Lemma 1, Proposition 1 in order of appearance, whereas the article prints the same items with the numbering of its own full document.  The complete native source of the article as distributed in the arXiv e-print for this exact version is bundled unchanged as \texttt{sources/201116/source0.tex}.
    \begin{lemma}\label{automorphism_characterization_of_D_l}
        If $l \geq 3$, then $\mathrm{Aut}(D_{l}) = \{f_{a,b} : 0 \leq a \leq l - 1, 1 \leq b \leq l - 1, \gcd(b, l) = 1\}$ where $f_{a,b}: D_{l} \to D_{l}$ is the group homomorphism defined by $f_{a, b}(r) = r^b$, $f_{a,b}(s) = r^as$.
    \end{lemma}

    \begin{proposition}\label{gcd_main_supporting_result}
        Let $x_1, x_2, x_3$ be positive integers. Then $\gcd (x_1, x_2) + \gcd (x_3, x_2) \leq x_2 + \gcd (x_1,x_3)$.
    \end{proposition}

    \begin{proposition}\label{gcd_secondary_supporting_result}
        Let $x_1$ be a natural number. Then
        \begin{align*}
            \gcd(x_1-1,x_1+1)=
            \begin{cases}
                2, & \text{if $x_1$ is odd},\\[4pt]
                1, & \text{if $x_1$ is even}.
            \end{cases}
        \end{align*}
    \end{proposition}

       \begin{lemma}\label{identity_greater_other}
            Let $\sigma$ be an automorphism of $D_l$ and $e$ denote the identity automorphism of $D_l$. Then $m_e \geq m_{\sigma}$.
        \end{lemma}

        \begin{proof}
            
             Since reflection is of the form \(r^ks \), where \(0\le k<l\). On employing Lemma \ref{automorphism_characterization_of_D_l}, we have
             \[
                \varphi_{a,b}(r^ks)=\varphi(r^k)\varphi(s) = r^{ak}\cdot r^bs  = r^{ak+b}s.
             \]
            Thus, \(r^ks\) is an involution iff \(r^{ak+b}s =  r^ks\), i.e.,
            \[
                r^{ak+b} = r^k \quad\Longleftrightarrow\quad (a-1)k \equiv -b \pmod l. \tag{1}
            \]
            Suppose $d_1=\gcd(a-1,l)$, the congruence \((a-1)k\equiv -b\pmod l\) has a solution iff \(d_1\mid b\), and has exactly \(d_1\) incongruent solutions modulo \(l\). So that, the number of reflections that are involutions equals \(d_1\) when \(d_1\mid b\), and equals \(0\) otherwise.
            And every rotation is of the form \(r^k\), where \(0\le k<l\). Then $\varphi_{a,b}(r^k)=\varphi(r)^k=r^{ak}$ and \(r^k\) is an involution iff \(r^{ak} = r^{-k}\), i.e.,
            \[
                r^{ak} = r^{-k} \quad\Longleftrightarrow\quad (a+1)k \equiv 0 \pmod l. \tag{2}
            \]
            Now, let \(d_2=\gcd(a+1,l)\), then the congruence \((a+1)k\equiv 0\pmod l\) has a solution, because \(d_2\mid 0\), and has exactly \(d_2\) incongruent solutions modulo \(l\). So that, the number of rotations that are involutions are $d_2$. Now, on summing up the maximum number of reflections that are involutions of $\varphi_{a, b}$ and maximum number of rotations that are involutions of $\varphi_{a, b}$, we obtain $\gcd(a - 1, l) + \gcd(a + 1, l)$. Therefore, the result follows from Lemma \ref{gcd_main_supporting_result}, and Lemma \ref{gcd_secondary_supporting_result}. This completes our argument.
        \end{proof}

\end{document}
